% Practical Exercise No. 7 about Gamma Radiation Attenuation

\section*{ATENUACIÓN DE LA RADIACIÓN GAMMA}

\textbf{OBJETIVO:} Después de esta práctica, los alumnos deben saber ampliar y determinar cómo pueden protegerse de las radiaciones utilizando para ello los diferentes blindajes existentes.

\subsection*{MATERIAL}
\begin{enumerate}
    \item Equipo Portátil Detector de Radiación tipo Geiger-Müller.
    \item Contenedor de Gammagrafía. (Conteniendo una fuente radiactiva con la suficiente actividad como para poder tener en sus costados niveles de rapidez de exposición de alrededor de 50 a 100 mR/h).
    \item Láminas de Hierro delgadas de 30 X 30 cm.
\end{enumerate}

\subsection*{DESARROLLO}
Se coloca el contenedor en el suelo de tal manera que uno de sus costados quede viendo hacia arriba, se miden las lecturas de rapidez de exposición a contacto de ese mismo costado del contenedor. Se anota el dato en una libreta de control. Se coloca una de las láminas metálicas sobre el costado del equipo y se miden nuevamente las lecturas de rapidez de exposición; esta misma acción se repite con cada una de las láminas (segunda y tercera), registrando los datos obtenidos en la libreta mencionada. Con estos datos se construye una gráfica como la del ejemplo siguiente: 

% Following this, a graphical example would typically be presented to illustrate the results.
